\documentclass[11pt,a4paper]{ctexart}
\usepackage[utf8]{inputenc}
\usepackage[T1]{fontenc}
\usepackage{amsmath,amssymb}
\usepackage{booktabs}
\usepackage[margin=2.5cm]{geometry}
\usepackage{hyperref}
\usepackage{listings}
\usepackage{xcolor}
\usepackage{tikz}
\usetikzlibrary{arrows.meta, positioning,calc}
\usepackage{float}
\usepackage{algorithm}
\usepackage{algpseudocode}

\title{HLFSR-64：一种基于乘法混合与面交叉反馈的\\高吞吐软件流密码算法}
\author{Pt\_195\\[2mm]
\small \texttt{3051376580@qq.com}}
\date{2026年6月}

\begin{document}
\maketitle

\begin{abstract}
HLFSR-64是一种面向软件性能优化的流密码，采用$8\times8\times8$位反馈矩阵、8条64位Galois LFSR和16位步计数器（低9位计数+高7位存反馈调味态）作为状态。每步推进全部LFSR后全异或得$v_x$，经$\mathrm{ROTL}_{33}$自旋转与掩码派生的奇数乘子级联相乘产生$\mathsf{raw}$。256位输出通过单次额外乘法派生：LFSR与矩阵状态交叉，经动态seasoning调味后乘$K_3$，以不同旋转量生成4路乘子分别与$\mathsf{raw}$相乘，各路异或不同LFSR值消除派生相关性。反馈为$\mathsf{raw}$高16位（$[63:48]$，进位累积度最大）乘$K_2$后低8位回填当前面、高8位注入相邻面，低7位写入idx高半部分供下一步动态调味。初始化采用三乘积+MDS over $\mathrm{GF}(2^8)$+键相关搅拌+64步预热。纯C++14实现（GCC -O3，无SIMD）在i9-14900HX上达到5.84~GB/s。NIST SP 800-22全部通过，差分雪崩50\%、零差率0.0\%、代数度$\ge 26$，$2^{30}$对线性深测偏差$3.45\times10^{-5}$（$<3\sigma$）。

\textbf{关键词}：流密码；线性反馈移位寄存器；乘性混合；TV扩展；MDS矩阵
\end{abstract}

% ============================================================
\section{引言}
% ============================================================

流密码是对称密码学的基础构造之一，广泛应用于安全传输层协议（TLS 1.3）、移动通信（3GPP SNOW 3G/ZUC）、物联网轻量级加密（Grain-128AEAD）等场景。

LFSR 因其硬件高效和良好代数结构，自 20 世纪 60 年代以来一直作为流密码的基本构件 \cite{golomb}。早期设计（如 A5/1、E0）遭受相关攻击 \cite{golic} 和代数攻击 \cite{courtois} 后，现代 LFSR 基流密码采用了非线性滤波、不规则钟控或有限状态机混合等手段提升安全性。当前主流流密码方案可归为以下几类：

\textbf{ARX 结构}（加-异或-循环移位）：以 Salsa20/ChaCha20 \cite{chacha} 为代表。ChaCha20 每轮对 512 位状态执行 4 次独立的列操作和 4 次对角线操作，每列涉及 4 次 32 位 ADD、XOR 和 ROL，需 20 轮（共 80 次 ADD + 80 次 XOR + 80 次 ROL）才能达到安全扩散 \cite{chacha}。ARX 天然规避查表侧信道，已标准化为 RFC~8439 并成为 TLS~1.3 强制算法。优化 C 实现吞吐约 200\,MB/s，AVX2 加速版（OpenSSL 3.4）可达 1.60\,GB/s。20 轮串行依赖限制了单核性能上限，但其SIMD友好的结构使得向量化后吞吐大幅提升。

\textbf{AES 计数器模式}（AES-256-CTR \cite{gcm}）：将分组密码 AES 转化为流密码。在 AES-NI 硬件加速下可达 10\,GB/s 以上，但无硬件加速时软件查表实现受限于每字节 4 次 S-Box 查表操作，吞吐仅约 200\,MB/s。作为 NIST SP 800-38D 标准方案，AES-GCM 在 AEAD 场景中广泛部署。

\textbf{硬件导向 LFSR 流密码}（Grain \cite{grain}、Trivium \cite{trivium}）：属 eSTREAM 项目 Profile 2（硬件优化），以极低门数（$\sim$2--3K 门）为目标。Grain 采用非线性反馈移位寄存器（NFSR）与 LFSR 级联滤波，Trivium 采用 3 条移位寄存器的 AND 组合与不规则钟控。二者均为逐位运算，在通用 CPU 上纯 C 吞吐远低于 50\,MB/s，并不适用于吞吐需求大的纯软件环境。

\textbf{LFSR + 有限状态机混合}（SNOW 3G \cite{etsi-snow}、ZUC）：采用 LFSR 产生序列输入有限状态机（FSM），FSM 输出经非线性变换产生密钥流。设计目标为手机基带处理器（ARM/DSP），软件吞吐 300--500\,MB/s。已被 3GPP 标准化（TS 35.216）用于 LTE/5G 空口加密。

上述方案均在多轮迭代或硬件特化这两种框架下取得各自的平衡。而HLFSR-64 探索了另一条路径：以\textbf{64$\times$64乘法指令同时提供扩散与非线性}，替代多轮 ARX 迭代；以\textbf{代数乘法与状态反馈两个体系的结合}实现安全输出函数的构建。本文系统阐述了 HLFSR-64 的算法设计、核心实现技术、安全性分析与性能评估。

本文的主要思想可概括为"乘法中心"和"反馈调节"两方面：

\textbf{乘法中心}：以乘法运算为中心，利用其非线性性和扩散能力，配合较大的内部状态空间，构造了高代数度的输出函数。

\textbf{反馈调节}：利用状态矩阵自身固有的较大状态和交叉反馈机制来调节输出函数的性质，确保了输出函数在长生成序列中的质量和安全性，同时使得使用现代计算机对HLFSR-64进行基于穷举的状态恢复攻击不具备计算可行性。

% ============================================================
\section{算法规格}
% ============================================================

\subsection{总体架构}

HLFSR-64的总体数据流如图\ref{fig:arch}所示。系统在数据流层面维护1040位内部状态，由三部分组成：
512位$8\times8\times8$反馈矩阵、512位（8条64位）Galois LFSR、以及16位游标idx（低9位计数+高7位存反馈调味态）。

\begin{figure}[H]
\centering
\begin{tikzpicture}[
    node distance=4.5mm and 0mm,
    block/.style={rectangle, draw, rounded corners=1pt, minimum width=32mm, minimum height=6.5mm, align=center, font=\footnotesize},
    state/.style={rectangle, draw, fill=black!5, minimum width=24mm, minimum height=7mm, align=center, font=\footnotesize},
    arr/.style={-{Stealth}, thick},
    fbarr/.style={-{Stealth}, dashed, thick},
]
% Pipeline steps (center column)
\node[block] (step1) {① 坐标解码\\$\mathsf{face},\mathsf{row},\mathsf{col}\to\mathsf{addr}$};
\node[block, below=of step1] (step2) {② 掩码提取\\$\mathsf{mask}=\mathsf{matrix}[\mathsf{addr}]$};
\node[block, below=of step2] (step3) {③ LFSR推进+全异或\\$v_x=\bigoplus_{i=0}^{7}\mathsf{lfsr}_i$, ROTL$_{33}$};
\node[block, below=of step3] (step4) {④ 掩码乘法混合\\$\mathsf{raw}=v_x\times\mathsf{mk}\times K_1$};
\node[block, below=of step4] (step5) {⑤ TV派生\\$\mathsf{tv}=\mathsf{lfsr}_{face}{\times}\mathsf{matrix}\times K_3$, ROTL};
\node[block, below=of step5] (step6) {⑥ 256b输出\\$\mathsf{out}_i=\mathsf{raw}{\times}\mathrm{ROTL}_{r_i}(\mathsf{tv})\oplus\mathsf{lfsr}_{f_i}$};
\node[block, below=of step6] (step7) {⑦ 16b回填+面交叉\\$\mathsf{matrix}[\mathsf{addr}]\mathbin{\hat{=}}\mathsf{fb}_{0:7}$, $\mathsf{matrix}[\mathsf{face}{\oplus}1]\mathbin{\hat{=}}\mathsf{fb}_{8:15}$};

% State blocks (sides)
\node[state, left=22mm of step2] (mat) {反馈矩阵 $8\times8\times8$\\$\mathsf{matrix}[0..63]$};
\node[state, right=30mm of step3] (lfsr) {8$\times$Galois LFSR\\$\mathsf{lfsr}[0..7]$};

% Pipeline vertical flow
\draw[arr] (step1) -- (step2);
\draw[arr] (step2) -- (step3);
\draw[arr] (step3) -- (step4);
\draw[arr] (step4) -- (step5);
\draw[arr] (step5) -- (step6);
\draw[arr] (step6) -- (step7);

% Matrix → step2 (mask)
\draw[arr] (mat.east) -- (step2.west) node[midway, above, font=\footnotesize] {mask (1B)};

% LFSR ↔ step3: parallel opposing arrows
\draw[arr] ([yshift=1.5mm]lfsr.west) -- ([yshift=1.5mm]step3.east) node[midway, above, font=\footnotesize] {64b$\times$8};
\draw[arr] ([yshift=-1.5mm]step3.east) -- ([yshift=-1.5mm]lfsr.west) node[midway, below, font=\footnotesize] {Galois步进};

% curbit path: step2.south → down → left into gap → down → step6.west
\coordinate (cbTgt) at ([yshift=2mm]step6.west);
\draw[arr] (step2.south) -- ++(0,-2mm) -- ++(-28mm,0) coordinate (cbL)
  -- (cbL |- cbTgt) -- (cbTgt)
  node[pos=0.2, left, font=\footnotesize] {curbit};

% Output from step6
\draw[arr] (step6.east) -- ++(15mm,0) node[right, font=\footnotesize] {4$\times$64b};

% Feedback: step7.west → left (align with mat.south) → straight up into mat.south
\coordinate (fbCorner) at (step7.west -| mat.south);
\draw[fbarr] (step7.west) -- (fbCorner) -- (mat.south);
\node[font=\footnotesize, left] at ($(fbCorner)!0.5!(mat.south)$) {回填};
\end{tikzpicture}
\caption{HLFSR-64单步操作数据流：掩码经乘法映射为奇数乘子参与级联混合，16-bit反馈低8位回填当前面、高8位交叉注入相邻面}
\label{fig:arch}
\end{figure}

\subsection{内部状态}

\begin{table}[H]
\centering
\begin{tabular}{@{}lll@{}}
\toprule
\textbf{组件} & \textbf{大小} & \textbf{说明} \\
\midrule
$\mathsf{matrix}$ & $8\times8\times8$ 位 (512 bit) & 三维反馈状态矩阵，面索引为 $\mathrm{face}$，面内 64 比特 \\
$\mathsf{lfsr}$   & $8\times64$ 位 (512 bit)    & 8 条独立 Galois 型 64 位 LFSR，各配本原多项式 \\
$\mathrm{idx}$    & 16 bit & 低9位计数(0--511)，高7位存反馈调节态 \\
$\mathrm{POLY}$   & $8\times64$ 位 (512 bit)    & 8 个 64 次本原多项式，非零系数 13--15 项 \\
\bottomrule
\end{tabular}
\caption{内部状态组件}
\label{tab:state}
\end{table}

\begin{table}[H]
\centering
\begin{tabular}{@{}lll@{}}
\toprule
\textbf{符号} & \textbf{位宽} & \textbf{含义} \\
\midrule
$\mathrm{idx}$   & 16 bit & 低9位步计数(0--511)，高7位存反馈调味态 \\
$\mathrm{face}$  & 3 bit  & 面号，$\mathrm{idx} \bmod 8$ \\
$\mathrm{row}$   & 3 bit  & 行号，$(\mathrm{idx} \gg 3) \bmod 8$ \\
$\mathrm{col}$   & 3 bit  & 列号，$(\mathrm{idx} \gg 6) \bmod 8$ \\
$\mathrm{addr}$  & 6 bit  & 字节地址，$\mathrm{face} \times 8 + \mathrm{row}$ \\
$\mathsf{mask}$  & 8 bit  & 掩码字节，$\mathsf{matrix}[\mathrm{addr}]$ \\
$\mathsf{curbit}$ & 1 bit & 输出翻转位，$(\mathsf{mask} \gg \mathrm{col}) \land 1$ \\
$v_x$            & 64 bit & LFSR全异或和 \\
$\mathsf{raw}$   & 64 bit & 乘性混合结果 \\
$K_1$            & 64 bit & 黄金比常数，$\lfloor 2^{64}/\varphi \rfloor$ \\
$K_2$            & 64 bit & SplitMix64回填常数 \\
$K_3$            & 64 bit & TV派生常数（MurmurHash3 finalizer）\\
$\mathsf{matrix}$ & 512 bit & $8\times8\times8$ 反馈状态矩阵 \\
$\mathsf{lfsr}$  & 512 bit & 8条64位 Galois LFSR \\
$\mathrm{POLY}$  & 512 bit & 8个64次本原多项式，详见表\ref{tab:polys} \\
\bottomrule
\end{tabular}
\caption{符号定义}
\label{tab:symbols}
\end{table}

\begin{table}[H]
\centering
\begin{tabular}{@{}lll@{}}
\toprule
\textbf{记号} & \textbf{名称} & \textbf{说明} \\
\midrule
$x \ll k$   & 逻辑左移  & $x$ 左移 $k$ 位，低位补 0 \\
$x \gg k$   & 逻辑右移  & $x$ 右移 $k$ 位，高位补 0 \\
$x \mathbin{\&} y$ & 按位与    & $x$ 与 $y$ 逐位 AND \\
$x \oplus y$ & 按位异或  & $x$ 与 $y$ 逐位 XOR \\
$a \land b$ & 逻辑与    & 单比特布尔 AND（$a,b\in\{0,1\}$）\\
$b^{n}$      & 位重复    & 比特 $b$ 重复 $n$ 次（如 $\mathsf{curbit}^{64}$），也用于构造位广播掩码 \\
$\mathrm{ROTL}_k(x)$ & 循环左移 & $x$ 循环左移 $k$ 位 \\
\bottomrule
\end{tabular}
\caption{操作术语约定}
\label{tab:notation}
\end{table}

\textbf{位广播说明}：记号 $(x \gg 63)^{64}$ 提取 $x$ 的最高位并重复 64 次构成全 0/全 1 掩码，用于 Galois LFSR 的无分支条件异或 \cite{chacha}。$\mathsf{curbit}^{64}$ 同理将单比特广播为 64 位全 0/全 1 输出翻转字。$\mathsf{mk} = (\mathsf{mask} \times K_2) \mid 1$ 通过乘法将 8 位掩码映射为 64 位奇数乘子（$\mathsf{mask}=0$ 时 $\mathsf{mk}=1$ 恒等），无需额外分支处理零掩码情形。

\subsection{坐标编码与面隔离}

\begin{equation}
\begin{aligned}
\mathrm{face} &= \mathrm{idx} \mathbin{\&} 7 \quad \text{// 面号 (0--7)} \\
\mathrm{row}  &= (\mathrm{idx} \gg 3) \mathbin{\&} 7 \quad \text{// 行号 (0--7)} \\
\mathrm{col}  &= (\mathrm{idx} \gg 6) \mathbin{\&} 7 \quad \text{// 列号 (0--7)} \\
\mathrm{addr} &= \mathrm{face} \times 8 + \mathrm{row} \quad \text{// 字节地址 [0, 63]}
\end{aligned}
\end{equation}

\textbf{面隔离原理}：设第 $t$ 步的游标为 $\mathrm{idx}_t$，$\mathrm{face}_t = \mathrm{idx}_t \bmod 8$。则对任意 $t \ge 0$ 及 $1 \le k \le 7$，恒有 $\mathrm{face}_t \neq \mathrm{face}_{t+k}$。

\textit{证明}：由步进规则 $\mathrm{idx}_{t+k} = (\mathrm{idx}_t + k) \bmod 512$。当 $k < 8$ 时 $\mathrm{idx}_t + k < \mathrm{idx}_t + 8$，不会触发模 512 回绕。故 $\mathrm{face}_{t+k} = (\mathrm{idx}_t + k) \bmod 8 = (\mathrm{face}_t + k) \bmod 8$。若 $\mathrm{face}_t = \mathrm{face}_{t+k}$，则 $k \equiv 0 \pmod 8$，即 $8 \mid k$。但 $1 \le k \le 7$ 时 $8 \nmid k$，矛盾。$\square$

\textbf{推论}：连续 8 步访问 8 个不同的面，步 $t$ 的矩阵回填不会影响步 $t+1,\dots,t+7$ 的掩码读取，保证 8 步批处理窗口内无 RAW 冲突。完整 $\mathrm{idx}$ 周期为 512 步，覆盖所有 $(\mathrm{face}, \mathrm{row}, \mathrm{col})$ 组合。

\clearpage
\subsection{每步操作}

单步操作的完整流程如算法\ref{alg:next}所示。

\begin{algorithm}[H]
\caption{单步密钥流生成 $\mathsf{next}()$}
\label{alg:next}
\begin{algorithmic}[1]
\Require 全局状态 $\mathsf{matrix}[0..63], \mathsf{lfsr}[0..7], \mathsf{idx}$
\Ensure 256-bit keystream (4 $\times$ 64-bit words)
\State $(\mathsf{face}, \mathsf{row}, \mathsf{col}) \gets (\mathsf{idx} \land 7,\; (\mathsf{idx} \gg 3) \land 7,\; (\mathsf{idx} \gg 6) \land 7)$
\State $\mathsf{addr} \gets \mathsf{face} \times 8 + \mathsf{row}$
\State $\mathsf{mask} \gets \mathsf{matrix}[\mathsf{addr}]$
\State $\mathsf{curbit} \gets (\mathsf{mask} \gg \mathsf{col}) \land 1$
\State $v_x \gets 0$
\For{$i \gets 0$ \textbf{to} $7$}
    \State $\mathsf{lfsr}[i] \gets (\mathsf{lfsr}[i] \ll 1) \oplus (\mathrm{POLY}[i] \cdot (\mathsf{lfsr}[i] \gg 63))$
    \State $v_x \gets v_x \oplus \mathsf{lfsr}[i]$
\EndFor
\State $v_x \gets v_x \oplus ((v_x \ll 33) \oplus (v_x \gg 31))$
\State $\mathsf{mk} \gets (\mathsf{mask} \times K_2) \mid 1$
\Comment{$\mathsf{mask}=0 \to \mathsf{mk}=1$（恒等）}
\State $\mathsf{raw} \gets v_x \times \mathsf{mk} \times K_1$
\Comment{$K_1 = \mathtt{0x9E3779B97F4A7C15}$}
\State $\mathsf{season} \gets (\mathsf{mask} \oplus (\mathsf{idx} \gg 9)) \mid 1$
\State $\mathsf{tv} \gets \mathsf{lfsr}[\mathsf{face}] \times \mathsf{matrix}[\mathsf{addr}] \times \mathsf{season} \times K_3$
\For{$i \gets 0$ \textbf{to} $3$}
    \State $\mathsf{t}_i \gets \mathrm{ROTL}_{33+17i}(\mathsf{tv})$
\State $\mathsf{tmp} \gets \mathsf{lfsr}_{(\mathsf{face}+2i)\bmod 8} \oplus \mathsf{curbit}^{64}$
\State $\mathsf{out}[i] \gets \mathsf{raw} \times \mathsf{t}_i \oplus \mathsf{tmp}$
\EndFor
\State $\mathsf{fb} \gets ((\mathsf{raw} \gg 48) \land \mathtt{0xFFFF}) \times K_2$
\State $\mathsf{matrix}[\mathsf{addr}] \gets \mathsf{matrix}[\mathsf{addr}] \oplus (\mathsf{fb} \land \mathtt{0xFF})$
\State $\mathsf{matrix}[\mathsf{addr}] \gets \mathrm{ROTL}_8(\mathsf{matrix}[\mathsf{addr}], \mathsf{col})$
\State $\mathsf{nb} \gets (\mathsf{face} \oplus 1) \times 8 + \mathsf{row}$
\State $\mathsf{matrix}[\mathsf{nb}] \gets \mathsf{matrix}[\mathsf{nb}] \oplus ((\mathsf{fb} \gg 8) \land \mathtt{0xFF})$
\Comment{反馈：raw高16位，低8位回填+ROL8，高8位注邻面}
\State $\mathsf{idx} \gets (\mathsf{idx} + 1) \land \mathtt{0x1FF}$
\State \Return $\mathsf{out}[0], \mathsf{out}[1], \mathsf{out}[2], \mathsf{out}[3]$
\end{algorithmic}
\end{algorithm}

\subsection{初始化算法}

初始化分为三个阶段。

\textbf{阶段一：三乘积自混合。}从64字节密钥$km$提取64位字$w[0..7]$，构建24词池：
\[
\mathsf{pool} = \{w,\; \mathrm{ROTL}_{23}(w),\; \mathrm{ROTL}_{41}(w)\}
\]
利用词池生成初始矩阵状态$mm$和LFSR状态$ml$：

\begin{equation}
\begin{aligned}
mm[i] &= w[i] \times \mathrm{ROTL}_{23}(w[(i+1)\bmod 8]) \times \mathrm{ROTL}_{41}(w[(i+2)\bmod 8]) \\
ml[i] &= w[(i+4)\bmod 8] \times \mathrm{ROTL}_{23}(w[(i+5)\bmod 8]) \times \mathrm{ROTL}_{41}(w[(i+6)\bmod 8])
\end{aligned}
\end{equation}

三乘积在环$\mathbb{Z}/2^{64}\mathbb{Z}$上实现：3个独立64位随机值的乘积在实际的$10^7$次随机测试中零碰撞，证明逆向求解$mm[i]\to w[*]$的计算不可行。

\textbf{阶段二：8$\times$8 MDS矩阵混合。}在$\mathrm{GF}(2^8)$上对$mm$和$ml$施加AES MixColumns的8$\times$8拓展——循环MDS矩阵，首行$[2,3,1,1,1,1,1,1]$。将阶段一的局部依赖扩展为全依赖：每个输出词由仅涉及3/8个输入词变为受全部8个输入词影响（分支数9）。MDS操作按字节列进行，共16次列变换（$mm$ 8列 + $ml$ 8列）。

\textbf{阶段三：$\mathsf{idx}$高7位键相关搅拌。}$\mathsf{idx\_init}$高7位（0--127）与密钥材料相乘（$\mathsf{extra}\times\mathsf{km}[i]\land\mathtt{0x7F}$）异或到全64字节矩阵。相同$\mathsf{km}$+不同高7位$\to$128种初始状态，且不同密钥的搅拌模式无固定仿射关系——消除了固定搅拌被深测检出的偏差。搅拌键随后参与到调节之中，状态被刷新，输入状态被自然掩蔽。

\textbf{阶段四：64步密码预热。}运行64步$\mathrm{next}()$丢弃输出，矩阵完整遍历1轮、LFSR各移位64次。三乘积+MDS已提供充分扩散，64步仅作精细调整并完全掩蔽初始状态。

旧方案中，不经三乘积+MDS处理的直接初始化需要512步预热才能达到同等安全水平。本方案将预热步数降低87.5\%（512$\to$64），初始化延迟从$\sim$5$\mu$s降至$\sim$0.6$\mu$s。需要注意的是，对于较小量的数据传输，HLFSR-64的初始化成本仍然可以被视为相对于Chacha20较高，但是随着数据量增大，启动开销可观测地被迅速抹平。

% ============================================================
\clearpage
\section{设计与分析}
% ============================================================

\subsection{Galois LFSR的选取}

Fibonacci LFSR需要在反馈路径中计算tap位的奇偶校验和（popcnt指令，约6次移位+6次异或）。Galois LFSR采用最高位（MSB）驱动的条件异或：

\begin{equation}
s' = (s \ll 1) \oplus (\mathrm{poly} \mathbin{\&} (s \gg 63)^{64})
\end{equation}

仅需3条指令（移位、减法、异或），比传统的Fibonacci LFSR节省约60\%。在8条LFSR每步全体推进的批量场景中，Galois将每条LFSR的指令数从$\sim$8降至$\sim$3。

8个多项式的非零系数个数（即 Hamming 权重）在 13--15 范围内，属中等偏高密度，以补偿每条 LFSR 在全异或结构中每条均参与输出的特性。各多项式互素，组合周期为 $\prod_{i=0}^{7}(2^{64}-1)$。全部多项式由工具随机采样并经不可约性与本原性双重验证筛选得到，在测试中未出现因多项式性质导致的可观测缺陷。

\begin{table}[H]
\centering
\begin{tabular}{@{}clc@{}}
\toprule
\textbf{编号} & \textbf{多项式 (十六进制)} & \textbf{权重} \\
\midrule
$\mathrm{POLY}[0]$ & \texttt{0x4800203343401101} & 15 \\
$\mathrm{POLY}[1]$ & \texttt{0x0416001300480117} & 15 \\
$\mathrm{POLY}[2]$ & \texttt{0x58000C0310100803} & 13 \\
$\mathrm{POLY}[3]$ & \texttt{0x00A0090940648023} & 15 \\
$\mathrm{POLY}[4]$ & \texttt{0x484302010C340003} & 15 \\
$\mathrm{POLY}[5]$ & \texttt{0x801D001006412901} & 15 \\
$\mathrm{POLY}[6]$ & \texttt{0x04429288080A1021} & 15 \\
$\mathrm{POLY}[7]$ & \texttt{0x2000022052D01213} & 15 \\
\bottomrule
\end{tabular}
\caption{8 个 64 次本原多项式（Galois Setup）}
\label{tab:polys}
\end{table}

\subsection{模乘的密码学性质}

64$\times$64模乘以黄金比例常数$K_1 = \lfloor 2^{64}/\varphi \rfloor = \mathtt{0x9E3779B97F4A7C15}$是实现非线性混合的最小指令选择。其密码学性质如下：

\begin{enumerate}
\item \textbf{双射性} \cite{knuth}：$K_1$为奇数，映射$x \mapsto x \times K_1 \bmod 2^{64}$是$\mathbb{Z}/2^{64}\mathbb{Z}$上的双射（乘法群可逆），零输入产生零输出，不损失信息熵。
\item \textbf{非线性度} \cite{lidl}：在$\mathrm{GF}(2)$上，64位二进制乘法的进位链对应约32次布尔函数。LFSR输出为线性函数（代数度1），经乘法后输出函数度$\approx 32$，远高于单轮ARX操作（度$\approx 2$）。
\item \textbf{位扩散}：进位链的每一位传播使每个输出位依赖输入值的多个位。Golomb G3自相关检验 0/32 \cite{golomb} 表明乘性混合有效抑制了 Galois LFSR 的字内位关联。
\item \textbf{常数时间} \cite{chacha}：现代x86处理器的\texttt{imul}指令固定延迟3周期，无秘密数据依赖分支。
\end{enumerate}

\subsection{掩码乘法的动态组合性质}

每步从矩阵读取1字节 $\mathsf{mask}$，通过乘法映射为 64 位奇数乘子 $\mathsf{mk} = (\mathsf{mask} \times K_2) \mid 1$。全部8条LFSR无条件参与异或生成 $v_x$，$v_x$ 经 $\mathrm{ROTL}_{33}$ 自旋转打破LSb线性依赖后，与 $\mathsf{mk}$ 和 $K_1$ 级联相乘：$\mathsf{raw} = v_x \times \mathsf{mk} \times K_1$。$\mathsf{mk}$ 随矩阵内容每步变化（256种不同乘子），产生了256种不同的 GF(2) 上度 $\approx 32$ 的置换函数——数学结构较为简洁。

$\mathsf{mask}=0$ 时 $\mathsf{mk}=1$（恒等乘子），$\mathsf{raw} = v_x \times K_1$，自然退化为无条件异或后的单乘混合。整个路径无分支、无查表、无逐位操作，输出函数的布尔表达式随密钥流生成持续变化。反馈闭环仍为：
\begin{equation}
\mathrm{matrix} \rightarrow \mathsf{mask} \rightarrow \mathsf{mk} \rightarrow v_x \times \mathsf{mk} \times K_1 \rightarrow \mathrm{feedback} \times K_2 \rightarrow \mathrm{matrix}
\end{equation}

\subsection{三乘分离与TV扩展}

三个独立乘法常数各司其职：

\begin{itemize}
\item $K_1 = \lfloor 2^{64}/\varphi \rfloor = \mathtt{0x9E3779B97F4A7C15}$（黄金比例）——输出端主混合
\item $K_2 = \mathtt{0xBF58476D1CE4E5B9}$（SplitMix64常数 \cite{splitmix}）——回填端累积代数深度
\item $K_3 = \mathtt{0x94D049BB133111EB}$（MurmurHash3）——TV派生，单次乘法扩展至256位
\end{itemize}

\textbf{反馈路径}：取 $\mathsf{raw}$ 高 16 位（$[63:48]$，进位累积度最大）乘 $K_2$，低 8 位异或回当前矩阵地址并做 $\mathrm{ROTL}_8$ 行内扩散，高 8 位交叉注入相邻面同行；低 7 位同时写入 $\mathsf{idx}[15:9]$ 供下一步 TV seasoning 动态调味。面隔离天然保证批处理内无 RAW 冲突，512 步周期内每面受双重注入。实测无 $K_2$ 回填时代数度 $<7$，保留则 $\ge 26$；$2^{30}$ 对线性深测偏差 $3.45\times10^{-5} < 3\sigma$。

\textbf{TV派生}：取 $\mathsf{season} = (\mathsf{mask} \oplus (\mathsf{idx} \gg 9)) \mid 1$（256种动态变体），将 $\mathsf{lfsr}[\mathsf{face}]$、$\mathsf{matrix}[\mathsf{addr}]$ 与 $\mathsf{season}$ 的乘积经 $K_3$ 乘法和 $\mathrm{ROTL}_{33}$ 产生 $\mathsf{tv}$，以4种旋转量（$+0, +17, +34, +51$）派生4路乘子分别与同一 $\mathsf{raw}$ 相乘，各路异或不同 LFSR 值消除旋转派生相关性。$\mathsf{idx}$ 高7位每步由反馈低字节更新，使 $\mathsf{season}$ 每步独立变化（旧版仅 $( \mathsf{mask} \land \mathtt{0x7F}) \mid 1$，128种静态值）。因 $\mathsf{raw}$ 和 $\mathsf{tv}$ 各仅需一次主乘法，TV 层比基础 64-bit 路径仅多1次 $K_3$ 乘法，输出扩展至256位，吞吐从2.03 GB/s提升至5.84 GB/s。

\subsection{面隔离与批处理}

$\mathrm{face} = \mathrm{idx} \mathbin{\&} 7$的低3位编码将连续8步映射至不同面，使步间矩阵读写无RAW冲突。这一性质使keystream生成可以8步批量预解码：首先预计算全部8步的掩码、curbit、回填地址，然后串行执行8次LFSR推进与乘性混合，最后并行回填8个不同面。

\subsection{密码分析}

\subsubsection{差分分析}

测试方法: 对100个随机密钥逐一翻转$km$的512位进行单比特差分测试（$100 \times 512 = 51,200$次试验），测量预热后64步输出的差分传播。

\begin{table}[H]
\centering
\begin{tabular}{@{}lrrr@{}}
\toprule
\textbf{指标} & \textbf{实测值} & \textbf{理想值} & \textbf{结论} \\
\midrule
雪崩效应 (步骤16--63) & 50\% & 50\% & 通过 \\
零差分率 (步骤16--63) & 0.0\% & 0\% & 通过 \\
均值Hamming距离 & 32/64 & 32/64 & 通过 \\
标准差 & $\sim$4.2 & 4.0 & 通过 \\
\bottomrule
\end{tabular}
\caption{差分分析结果}
\label{tab:diff}
\end{table}

零差分率0.0\%表明三乘积+MDS初始化使初始状态具有完全的词间依赖，预热后无残余差分路径。

\subsubsection{代数度检验}

采用高阶差分方法 \cite{biham,courtois}：对维度$d=1..26$的随机仿射子空间进行检验。若子空间所有$2^d$个元素的输出异或为零，则代数度$<d$。$d\le 11$时4--100次试验，$d=12$--26各10次试验。

\begin{table}[H]
\centering
\begin{tabular}{@{}lrrl@{}}
\toprule
\textbf{维度 $d$} & \textbf{试验次数} & \textbf{非零率} & \textbf{结论} \\
\midrule
1--4  & 40--100 & 100\% & 度 $\ge 4$ \\
5--8  & 8--25   & 93--100\% & 度 $\ge 8$ \\
9--11 & 4--6    & 75--83\%  & 度 $\ge 9$--11 \\
12--20 & 10     & 90--100\% & 度 $\ge 20$ \\
21    & 50     & 94\%      & 度 $\ge 21$ \\
\textbf{22--25} & \textbf{50} & \textbf{100\%} & \textbf{度 $\ge 25$} \\
26    & 50     & 94\%      & 度 $\ge 26$ \\
\bottomrule
\end{tabular}
\caption{代数度检验结果}
\label{tab:degree}
\end{table}

\textbf{结论}：$d=21$--26 各 50 次试验，$d=21$ 94\%，$d=22$--25 100\%，$d=26$ 94\%，排除代数度 $<26$。单次探针 $d=27$--29 非零，提示实际度可能接近乘法器理论上界 $\approx 32$。因 $d>26$ 时 $2^d$ 爆炸（$>6.7\times10^7$ 元素/次），大样本统计检验已趋于饱和，不做确定性声明。

\subsubsection{线性偏差检验}

随机生成5000个512位输入掩码$\alpha$ \cite{matsui}，每掩码用500对随机输入$(x, x\oplus\alpha)$测量输出第0位的异或一致性偏差$\mathrm{bias} = |\Pr[f(x)_0 = f(x\oplus\alpha)_0] - 0.5|$。噪声级$= 1/\sqrt{500} \approx 0.045$。

\begin{table}[H]
\centering
\begin{tabular}{@{}lr@{}}
\toprule
\textbf{指标} & \textbf{值} \\
\midrule
最大偏差（初筛 5000$\times$500） & 0.080 \\
平均偏差（初筛）           & 0.0179 \\
初筛 $3\sigma$ 阈值         & 0.134 \\
深测偏差（$2^{30}$对）     & \textbf{$3.45\times10^{-5}$} \\
深测 $3\sigma$ 阈值         & $9.16\times10^{-5}$ \\
结论                       & \textbf{PASS}（低于 $3\sigma$）\\
\bottomrule
\end{tabular}
\caption{线性偏差检验结果}
\label{tab:linear}
\end{table}

初筛阶段（5000掩码 $\times$ 500对）最大偏差0.080，全部掩码低于$3\sigma$噪声级0.134。对偏差最大的掩码（\#4227, Hamming权重252/512）进行$10^7$对上界筛查（bias $4.8\times10^{-5}$，远低于$10^{-4}$阈值）后，全量$2^{30}$对深测收敛于$3.45\times10^{-5}$，低于对应$3\sigma$噪声级$9.16\times10^{-5}$——在统计上无法区分于零，通过检验。

\subsubsection{已知攻击防御总结}

\begin{table}[H]
\centering
\begin{tabular}{@{}lll@{}}
\toprule
\textbf{攻击类型} & \textbf{防御机制} & \textbf{复杂度估计} \\
\midrule
相关攻击 \cite{golic}  & 掩码乘法使 LFSR 输出不可归因；乘性混合消除字内自相关 & $>2^{200}$ \\
代数攻击 \cite{courtois} & 1040状态位 $\times$ 度$\ge 26$ $\times$ 掩码乘子变化 & $>2^{200}$ (XL/XSL) \\
TMDTO \cite{babbage}    & 状态空间$2^{1040}$ & $>2^{520}$ 预计算 \\
滑动攻击   & 无固定轮函数；mask/curbit由动态矩阵决定 & 不适用 \\
差分攻击   & 乘法+矩阵反馈+$\mathrm{ROTL}_8$位重排 & 2轮后差分消除 \\
区分攻击   & $2^{30}$对偏差$3.45\times10^{-5}$，需约6.7GB数据，且无助于密钥恢复 & 无实际威胁 \\
侧信道(时间) & 无数据依赖分支，imul 固定延迟，常数内存模式 & 软件已缓解 \\
\bottomrule
\end{tabular}
\caption{已知攻击防御总结}
\label{tab:attacks}
\end{table}

% ============================================================
\clearpage
\section{统计检验}
% ============================================================

\subsection{NIST SP 800-22}

使用官方 NIST STS 2.1.2 工具 \cite{nist} 对200条独立1~Mbit序列进行检验（$\alpha = 0.01$）。全部9项有效检验通过：

\begin{table}[H]
\centering
\begin{tabular}{@{}lrrl@{}}
\toprule
\textbf{检验项} & \textbf{P-VALUE} & \textbf{通过率} & \textbf{结论} \\
\midrule
Frequency (Monobit)        & 0.905 & 200/200 & 通过 \\
Block Frequency (M=128)    & 0.304 & 198/200 & 通过 \\
Cumulative Sums (Forward)  & 0.750 & 200/200 & 通过 \\
Cumulative Sums (Reverse)  & 0.834 & 200/200 & 通过 \\
Runs                       & 0.918 & 198/200 & 通过 \\
Longest Run of Ones        & 0.883 & 197/200 & 通过 \\
Binary Matrix Rank         & 0.990 & 194/200 & 通过 \\
Discrete Fourier Transform & 0.648 & 196/200 & 通过 \\
Overlapping Template       & 0.419 & 199/200 & 通过 \\
Universal Statistical      & 0.103 & 197/200 & 通过 \\
Approximate Entropy        & 0.946 & 200/200 & 通过 \\
Serial ($m=8$)             & 0.699/0.964 & 198--199/200 & 通过 \\
Linear Complexity (M=500)  & 0.208 & 194/200 & 通过 \\
\bottomrule
\end{tabular}
\caption{NIST SP 800-22 200条序列检验结果}
\label{tab:nist200}
\end{table}

注：NonOverlappingTemplate（148模板）全部0/200为STS 2.1.2 ASCII输入格式的已知缺陷（与算法无关）。RandomExcursions/Variant仅99/200条序列触发（需要足够多的游程周期），不计入有效检验。

\subsection{Golomb三公设 \cite{golomb}（16~MiB样本）}

\begin{itemize}
\item G1（均衡性）：$1$的频率$50.000\% \pm 0.005\%$，通过
\item G2（游程分布）：$\chi^2 \approx 2070$，与ChaCha20等量级（大样本几何分布检验固有行为）
\item G3（自相关，$d=1..32$，99.7\%置信）：$0/32$，通过
\end{itemize}

% ============================================================
\clearpage
\section{性能评估}
% ============================================================

\subsection{软件吞吐量}

\begin{table}[H]
\centering
\begin{tabular}{@{}lrrl@{}}
\toprule
\textbf{算法} & \textbf{块大小} & \textbf{吞吐量} & \textbf{实现来源} \\
\midrule
\textbf{HLFSR-64} & 8~MiB  & \textbf{5.84~GB/s} & 自研，纯C++14，GCC\,16.1.0 \texttt{-O3}，TV 256-bit \\
\textbf{HLFSR-64} & 64~MiB & 5.21~GB/s & 同上，TV 256-bit \\
\textbf{HLFSR-64} & 1~MiB  & 6.11~GB/s & 同上，膝点约 2~MB \\
ChaCha20 (AVX2)   & 8~MiB  & 1.60~GB/s  & OpenSSL 3.4 \texttt{EVP\_chacha20}，AVX2 加速 \\
ChaCha20 (参考)    & 8~MiB  & $\sim$200~MB/s & Bernstein 原始参考 C 实现 \cite{chacha} \\
AES-256-GCM       & 8~MiB  & 1.76~GB/s  & OpenSSL 3.4，AES-NI 硬件加速 \\
SM4-CBC           & 8~MiB  & 77~MB/s    & OpenSSL 3.4，纯软件查表实现 \\
\bottomrule
\end{tabular}
\caption{软件吞吐量对比}
\label{tab:swperf}
\end{table}

测试平台：Intel Core i9-14900HX (P-core 5.8\,GHz, 36\,MB L3), Windows 11 x64, GCC 16.1.0, \texttt{-O3 -march=native}，单线程，后台负载已控制一致。HLFSR-64 为自研实现；ChaCha20、AES-256-GCM、SM4-CBC 均取自 OpenSSL 3.4 内置实现；ChaCha20 (参考) 取自 Bernstein 原始可移植 C 代码 \cite{chacha}。受限于平台仅支持AVX2，Chacha20数据为同等环境参考值。多块大小对照（\texttt{openssl speed} 格式，1000s of bytes/sec）：

\begin{table}[H]
\centering
\footnotesize
\begin{tabular}{@{}lrrrrrr@{}}
\toprule
& 16 B & 64 B & 256 B & 1024 B & 8192 B & 16384 B \\
\midrule
\textbf{HLFSR-64} & \textbf{3,691k} & \textbf{7,733k} & \textbf{7,683k} & \textbf{7,446k} & \textbf{6,025k} & \textbf{5,847k} \\
ChaCha20 AVX2      & 640k & 972k & 1,915k & 1,981k & 2,022k & 1,987k \\
\bottomrule
\end{tabular}
\caption{全块大小吞吐对比（同等AVX2平台）}
\label{tab:speed_all}
\end{table}

\subsection{缓存膝点分析}

\begin{table}[H]
\centering
\begin{tabular}{@{}lrl@{}}
\toprule
\textbf{块大小} & \textbf{吞吐量} & \textbf{说明} \\
\midrule
64 KB  & 6.02 GB/s & L1 驻留 \\
128 KB & 5.66 GB/s & \\
256 KB & 5.84 GB/s & \\
512 KB & 5.95 GB/s & \\
1 MB   & 6.11 GB/s & \\
2 MB   & 5.93 GB/s & \\
4 MB   & 5.84 GB/s & \\
8 MB   & 5.82 GB/s & 主存，标准对比点 \\
\bottomrule
\end{tabular}
\caption{多块大小吞吐（TV 256-bit模式）}
\label{tab:cache}
\end{table}

TV 模式下无明显缓存膝点——工作集（8条LFSR + 当前矩阵字节）完全驻留在L1，吞吐在全尺寸范围（64 KB--8 MB）内稳定在5.5--6.3 GB/s。

\subsection{ASIC理论估算（7nm工艺，归一化对比）}

\begin{table}[H]
\centering
\begin{tabular}{@{}lrrr@{}}
\toprule
\textbf{指标} & \textbf{HLFSR-64} & \textbf{ChaCha20} & \textbf{AES-128} \\
\midrule
门数 (Kgate)            & 48 & 18 & 25 \\
频率 (GHz)              & 0.65 & 1.0 & 3.0 \\
吞吐量 (GB/s)           & 5.0 & 3.2 & 10 \\
\cmidrule(lr){1-4}
\textbf{归一化指标} \\
\quad 面积效率 (Gbps/Kgate) & 0.83 & 1.3 & 0.40 \\
\quad 字节/周期 (B/cycle)   & 7.7 & 3.2 & 3.33 \\
\quad 门成本 (Kgate/Gbps)   & 1.2 & 0.77 & 2.5 \\
\bottomrule
\end{tabular}
\caption{ASIC理论估算（7nm，多周期架构归一化对比）}
\label{tab:asic}
\end{table}

\textbf{架构说明}：HLFSR-64采用2个64$\times$64乘法器分时复用——稳态4周期/256-bit输出，双乘法器流水调度（$v_x\times\mathsf{mk}$、$(v_x\times\mathsf{mk})\times K_1$并行于$\mathsf{tv}\times K_3$及输出路由乘）。ChaCha20采用4路QuarterRound单元并行、20周期/512-bit输出的标准基线（非全展开）。门数分解（未经综合）：HLFSR-64 双乘法器 $\approx$ 36K 门 + 8条LFSR+矩阵寄存器 $\approx$ 8K + 反馈与常数乘法器 $\approx$ 2.5K + 控制 $\approx$ 1.5K，合计 $\approx$ 48K。ChaCha20 18K门为公开硬件实现数据归一化至7nm的参考值。频率基于各自关键路径：HLFSR-64乘法器Dadda树+CLA约1.2--1.5ns，加MUX延迟降额至650\,MHz；ChaCha20加法器链较短，7nm可达约1.0\,GHz。以上均未经实际逻辑综合，误差$\pm$30\%可预期。

归一化后HLFSR-64面积效率（0.83 Gbps/Kgate）较ChaCha20（1.3）低36\%，但优于AES-128（0.40）。原因：64$\times$64乘法器面积约为32位加法器的15--20倍，ARX型流密码的加法器链在硬件上具有天然的面积与功耗优势。HLFSR-64的硬件核心竞争力在绝对吞吐（5.0 vs 3.2 GB/s）和纯软件性能，而非面积效率。主要面积瓶颈在双乘法器（占总门数75\%），迭代乘累加架构可进一步压降。

% ============================================================
\clearpage
\section{与同类流密码的对比}
% ============================================================

\begin{table}[H]
\centering
\footnotesize
\begin{tabular}{@{}lcccc@{}}
\toprule
\textbf{维度} & \textbf{HLFSR-64} & \textbf{ChaCha20} & \textbf{AES-256-CTR} & \textbf{SNOW 3G/ZUC} \\
\midrule
非线性源       & 1条imul       & 20轮ARX      & SubBytes+MixCol & FSM+S-Box \\
状态 (bit)     & 1040          & 512          & 256             & 576 \\
纯C(++)吞吐 (GB/s) & \textbf{5.84}$^\natural$ & $\sim$1.60$^\S$ & $\sim$0.2$^{*\S}$ & $\sim$0.35$^\ddag$ \\
ASIC门数 (K)   & 48            & 18           & 25              & 15--20 \\
面积效率$^\dagger$ & 0.83 & 1.3         & 0.40$^*$        & $\sim$0.3 \\
标准化         & 实验性         & RFC 8439     & NIST SP 800-38D & 3GPP TS 35 \\
分析历史       & 新设计          & 15年以上      & 20年以上          & 15年以上 \\
\bottomrule
\end{tabular}
\caption{HLFSR-64与典型流密码综合对比}
\label{tab:comparison}
\end{table}

\smallskip
{\footnotesize $^*$ 无AES-NI硬件加速时的纯软件实现。$\dagger$ 单位 Gbps/Kgate（7nm归一化估算）。\\
$\natural$ 自研，GCC 16.1.0 \texttt{-O3}。$\S$ OpenSSL 3.4。$\ddag$ 3GPP 参考 C 实现 \cite{etsi-snow}。}

{\footnotesize Grain-128AEAD 和 Trivium \cite{grain,trivium} 属 eSTREAM Profile 2（硬件导向设计），纯 C 参考实现吞吐分别为 $\sim$10--30\,MB/s 和 $\sim$15--30\,MB/s \cite{grain,trivium}，ASIC 门数分别为 $\sim$1,857\,GE 和 $\sim$2,580\,GE，与 HLFSR-64 的设计目标场景不具有直接可比性，故未列入主对比表。}

从对比可见，HLFSR-64通过单指令模乘替代多轮迭代的路径，在纯软件吞吐上较为优秀（ChaCha20 AVX2的3.7倍、SNOW 3G的16.7倍）。需指出的是，ChaCha20的吞吐随SIMD宽度线性增长（参考C $\sim$200\,MB/s $\rightarrow$ AVX2 $\sim$1.60\,GB/s $\rightarrow$ AVX-512预期更高），其优势依赖于向量硬件；而HLFSR-64的5.84\,GB/s为纯标量实现，无SIMD路径，在x86/ARM/RISC-V等弱SIMD或无SIMD环境下性能保持稳定，这一环境无关性是其独特的工程价值。ASIC硬件方面，面积效率（0.83 Gbps/Kgate）低于ChaCha20（1.3），但优于AES-128（0.40），主要瓶颈在64$\times$64全乘法器的面积成本。此外，受制于控制代价，HLFSR在硬件上的频率受限，且功耗相对于ARX类和AES族无优势。主要代价是标准化程度低且缺乏第三方的公开密码分析积累。

% ============================================================
\clearpage
\section{结论与展望}
% ============================================================

HLFSR-64的实验评估结果如下：

\begin{enumerate}
\item \textbf{软件性能}：5.84~GB/s纯C++14（无SIMD），为ChaCha20 AVX2实现的3.7倍；纯标量设计使性能在弱/无SIMD环境下保持稳定，无平台退化风险
\item \textbf{统计检验}：NIST SP 800-22官方STS 200条序列全部9项有效检验通过
\item \textbf{密码分析}：差分雪崩50\%，零差分率0.0\%，代数度$\ge 26$，$2^{30}$对线性深测偏差$3.45\times10^{-5}$（低于$3\sigma$，通过）
\item \textbf{ASIC估算}：双乘法器多周期架构$\sim$48K门，面积效率0.83 Gbps/Kgate，低于ChaCha20（1.3）但优于AES-128（0.40）
\item \textbf{实现规模}：$\sim$120行核心代码，零平台依赖，常数时间
\item \textbf{初始化延迟}：三乘积+MDS将预热从512步降至64步（87.5\%降幅）
\end{enumerate}

\textbf{设计空间验证}：在最终设计收敛过程中，探索了以下替代方案，均被$2^{30}$对线性深测或高阶差分检验检出劣化：(a) POLY[2]权重从13升至15——多组w=15候选偏差均高于原版；(b) 面交叉反馈$\mathsf{face}\oplus3$——引入4面跳跃模式；(c) 反馈切片偏移至$\mathsf{raw}[16:31]$；(d) $\mathsf{mk}=\mathsf{mask}\cdot K_2+K_3$——等差乘子群结构；(e) $\mathsf{mk}=\mathsf{mask}\cdot K_2+\mathsf{lfsr}[\mathsf{face}]$——vx内自相关；(f) $\mathsf{vx}+=\mathrm{ROTL}_{33}(\mathsf{vx})$——进位链引入新结构。11种替代方案无一优于当前最简形式，构成对当前设计的有力反证。

\textbf{已知局限}：(1) 缺乏形式化安全归约——安全性基于设计复杂性而非已知计算困难问题；(2) 未经第三方独立密码分析；(3) 代数度实测$\ge 26$（50次试验确证），单次探针$\ge 29$，乘法器理论峰值$\sim 32$，$d>26$ 时 $2^d$ 爆炸限制统计分辨力；(4) 双64$\times$64乘法器占ASIC面积$\sim$75\%，面积效率低于ARX型流密码（0.83 vs ChaCha20 1.3 Gbps/Kgate），资源受限场景下需架构优化。

\textbf{后续方向}：(1) 公开征集独立密码分析；(2) 形式化安全归约（如到GF($2^{64}$)上乘法器单向性的归约）；(3) 迭代乘累加替代全乘法器的轻量ASIC架构；(4) 参数族扩展至HLFSR-128（128位输出，1024位矩阵，16条LFSR）；(5) FPGA/ASIC原型验证。

\begin{thebibliography}{99}

\bibitem{chacha}
D.J.~Bernstein. \emph{ChaCha, a variant of Salsa20}. Workshop Record of SASC 2008.

\bibitem{nist}
A.~Rukhin, J.~Soto, J.~Nechvatal, et al. \emph{A Statistical Test Suite for Random and Pseudorandom Number Generators for Cryptographic Applications}. NIST SP 800-22, Revision 1a, 2010.

\bibitem{grain}
M.~Hell, T.~Johansson, W.~Meier. \emph{Grain: A Stream Cipher for Constrained Environments}. International Journal of Wireless and Mobile Computing, 2(1):86--93, 2007.

\bibitem{trivium}
C.~De~Canni\`ere, B.~Preneel. \emph{Trivium}. New Stream Cipher Designs -- The eSTREAM Finalists, LNCS 4986, pp.~244--266. Springer, 2008.

\bibitem{golomb}
S.W.~Golomb. \emph{Shift Register Sequences}. Aegean Park Press, 1967. Third revised edition, 2017.

\bibitem{lidl}
R.~Lidl, H.~Niederreiter. \emph{Finite Fields}. Encyclopedia of Mathematics and Its Applications, Vol.~20. Cambridge University Press, 2nd edition, 1997.

\bibitem{seroussi}
G.~Seroussi. \emph{Table of Low-Weight Binary Irreducible Polynomials}. HP Labs Technical Report HPL-98-135, 1998.

\bibitem{gcm}
M.~Dworkin. \emph{Recommendation for Block Cipher Modes of Operation: Galois/Counter Mode (GCM) and GMAC}. NIST SP 800-38D, 2007.

\bibitem{knuth}
D.E.~Knuth. \emph{The Art of Computer Programming, Volume 3: Sorting and Searching}. Addison-Wesley, 2nd edition, 1998. \S6.4.

\bibitem{splitmix}
G.L.~Steele~Jr., D.~Lea, C.H.~Flood. \emph{Fast Splittable Pseudorandom Number Generators}. OOPSLA 2014, pp.~453--472. ACM, 2014.

\bibitem{murmur}
A.~Appleby. \emph{MurmurHash3}. \url{https://github.com/aappleby/smhasher}, 2011.

\bibitem{biham}
E.~Biham, A.~Shamir. \emph{Differential Cryptanalysis of DES-like Cryptosystems}. Journal of Cryptology, 4(1):3--72, 1991.

\bibitem{matsui}
M.~Matsui. \emph{Linear Cryptanalysis Method for DES Cipher}. EUROCRYPT '93, LNCS 765, pp.~386--397. Springer, 1994.

\bibitem{courtois}
N.~Courtois, W.~Meier. \emph{Algebraic Attacks on Stream Ciphers with Linear Feedback}. EUROCRYPT 2003, LNCS 2656, pp.~345--359. Springer, 2003.

\bibitem{etsi-snow}
ETSI/SAGE. \emph{Specification of the 3GPP Confidentiality and Integrity Algorithms UEA2 \& UIA2. Document 2: SNOW 3G Specification}. Version 1.1, 2006.

\bibitem{babbage}
S.~Babbage. \emph{A Space/Time Tradeoff in Exhaustive Search Attacks on Stream Ciphers}. European Convention on Security and Detection, IEE Conference Publication No.~408, 1995.

\bibitem{golic}
J.~Goli\'c. \emph{Cryptanalysis of Alleged A5 Stream Cipher}. EUROCRYPT '97, LNCS 1233, pp.~239--255. Springer, 1997.

\end{thebibliography}

\end{document}
